
\documentclass{article}


\usepackage{lipsum}
\usepackage{lmodern}
\usepackage[T1]{fontenc}
\usepackage[utf8]{inputenc}
\usepackage{mathtools}
\usepackage{amssymb}
\usepackage{amsmath}
\usepackage{graphicx}
\usepackage{sectsty} % Para cambiar el tamanio de los titulos
\usepackage[lmargin=1in, rmargin=0.75in]{geometry}
\usepackage{xcolor}
%\usepackage{hyperref}
\usepackage[colorlinks = true,
linkcolor = black,
urlcolor  = blue,
citecolor = blue,
anchorcolor = blue]{hyperref}


\usepackage{tikz}
\usepackage{xcolor}
\usepackage{longtable}
\usepackage{ctable}
\newcolumntype{C}[1]{>{\centering\let\newline\\\arraybackslash\hspace{0pt}}m{#1}}
%TeXporter:DefineEmojiColors
\definecolor{c0a0808}{RGB}{10,8,8}
\definecolor{cff0000}{RGB}{255,0,0}
\definecolor{cff9700}{RGB}{255,151,0}
\definecolor{c1a0f00}{RGB}{26,15,0}
\definecolor{c782121}{RGB}{120,33,33}
\definecolor{cffffff}{RGB}{255,255,255}
\definecolor{cffefef}{RGB}{255,239,239}
\definecolor{c250303}{RGB}{37,3,3}
\definecolor{ce3dbdb}{RGB}{227,219,219}
\definecolor{cd45500}{RGB}{212,85,0}
\definecolor{cffd42a}{RGB}{255,212,42}
\definecolor{caa4400}{RGB}{170,68,0}
\definecolor{cff5555}{RGB}{255,85,85}
\definecolor{cffff00}{RGB}{255,255,0}
\definecolor{c5a5a5a}{RGB}{90,90,90}
\definecolor{ce9c900}{RGB}{233,201,0}
\definecolor{c2b0000}{RGB}{43,0,0}
\definecolor{c00ffff}{RGB}{0,255,255}
%TeXporter:EndDefineEmojiColors

%TeXporter:EmojiStressedFace
\def\CreateStressedFace{
	\begin{tikzpicture}[scale=0.5]
	%Head
	\pgfdeclareradialshading{ballshading}{\pgfpoint{-10bp}{10bp}}
	{color(0bp)=(yellow!40!white);
		color(9bp)=(yellow!55!white);
		color(18bp)=(yellow!60!orange);
		color(25bp)=(yellow!50!orange);
		color(50bp)=(black)}
	\pgfpathcircle{\pgfpoint{0cm}{0cm}}{2cm}
	\pgfshadepath{ballshading}{20}
	\pgfusepath{}
	\pgfdeclareradialshading{ballshading}{\pgfpoint{-10bp}{10bp}}
	{color(0bp)=(yellow!40!white);
		color(9bp)=(yellow!30!white);
		color(18bp)=(yellow!60!orange);
		color(25bp)=(yellow!50!orange);
		color(50bp)=(orange)}
	\pgfpathcircle{\pgfpoint{0cm}{0cm}}{1.85cm}
	\pgfshadepath{ballshading}{20}
	\pgfusepath{}
	%Head hat
	\path[yshift=138,xshift=-140,scale=0.012,yscale=-1, fill=white,opacity=0.234] (411.4196,246.6473) .. controls
	(344.7661,246.6473) and (288.7870,289.5501) .. (273.0134,347.5848) .. controls
	(311.3176,336.0898) and (355.6292,329.4911) .. (402.8571,329.4911) .. controls
	(457.9662,329.4911) and (509.1043,338.4615) .. (551.3571,353.7411) .. controls
	(537.9844,292.6125) and (480.4076,246.6473) .. (411.4196,246.6473) -- cycle;
	
	\path[xshift=-165,yshift=190,scale=0.013,yscale=-1,draw=brown!70!black,inner color=c782121!70!yellow,outer color=c782121!60!black]
	(412.2952,416.1696) .. controls (402.4957,434.4069) and (372.7147,447.8466) ..
	(337.0714,448.6326) .. controls (345.7428,453.7104) and (357.9211,456.8754) ..
	(371.4088,456.8754) .. controls (397.7352,456.8754) and (419.0639,444.8410) ..
	(419.0639,429.9941) .. controls (419.0639,424.9402) and (416.5958,420.2082) ..
	(412.2952,416.1696) -- cycle;
	
	\path[xshift=165,yshift=190,scale=0.013,yscale=-1,xscale=-1,draw=brown!70!black,inner color=c782121!70!yellow,outer color=c782121!60!black]
	(412.2952,416.1696) .. controls (402.4957,434.4069) and (372.7147,447.8466) ..
	(337.0714,448.6326) .. controls (345.7428,453.7104) and (357.9211,456.8754) ..
	(371.4088,456.8754) .. controls (397.7352,456.8754) and (419.0639,444.8410) ..
	(419.0639,429.9941) .. controls (419.0639,424.9402) and (416.5958,420.2082) ..
	(412.2952,416.1696) -- cycle;
	
	\path[xshift=-51,yshift=28,scale=0.003,yscale=-1,xscale=1,draw=brown!70!black,inner color=c782121!70!yellow,outer color=c782121!60!black]
	(240.9449,201.9685) .. controls (277.7444,163.7790) and (331.4846,142.5483) ..
	(384.4488,145.2756) .. controls (416.1699,146.9090) and (447.4278,156.8390) ..
	(474.2758,173.8119) .. controls (501.1239,190.7848) and (523.5012,214.7621) ..
	(538.5827,242.7165) .. controls (491.0752,298.2483) and (428.6499,340.8953) ..
	(359.6457,364.9606) .. controls (319.2364,379.0534) and (276.6393,386.8529) ..
	(233.8583,387.9921) .. controls (271.0447,373.5979) and (305.8888,353.1720) ..
	(336.6142,327.7559) .. controls (350.5885,316.1962) and (363.9856,303.2534) ..
	(372.0472,287.0079) .. controls (380.5313,269.9111) and (382.4928,249.6748) ..
	(377.4504,231.2669) .. controls (372.4079,212.8589) and (360.4075,196.4472) ..
	(344.3949,186.0607) .. controls (328.3824,175.6743) and (308.5040,171.4077) ..
	(289.6397,174.3085) .. controls (270.7753,177.2094) and (253.0971,187.2510) ..
	(240.9449,201.9685);
	
	\path[xshift=51,yshift=28,scale=0.003,yscale=-1,xscale=-1,draw=brown!70!black,inner color=c782121!70!yellow,outer color=c782121!60!black]
	(240.9449,201.9685) .. controls (277.7444,163.7790) and (331.4846,142.5483) ..
	(384.4488,145.2756) .. controls (416.1699,146.9090) and (447.4278,156.8390) ..
	(474.2758,173.8119) .. controls (501.1239,190.7848) and (523.5012,214.7621) ..
	(538.5827,242.7165) .. controls (491.0752,298.2483) and (428.6499,340.8953) ..
	(359.6457,364.9606) .. controls (319.2364,379.0534) and (276.6393,386.8529) ..
	(233.8583,387.9921) .. controls (271.0447,373.5979) and (305.8888,353.1720) ..
	(336.6142,327.7559) .. controls (350.5885,316.1962) and (363.9856,303.2534) ..
	(372.0472,287.0079) .. controls (380.5313,269.9111) and (382.4928,249.6748) ..
	(377.4504,231.2669) .. controls (372.4079,212.8589) and (360.4075,196.4472) ..
	(344.3949,186.0607) .. controls (328.3824,175.6743) and (308.5040,171.4077) ..
	(289.6397,174.3085) .. controls (270.7753,177.2094) and (253.0971,187.2510) ..
	(240.9449,201.9685);
	
	\path[xshift=66,yshift=1,scale=0.0065,yscale=-1,xscale=-1,draw=black,inner color=cff5555!80!black,outer color=c782121!60!black]
	(359.6464,67.3100) .. controls (278.1326,67.3100)
	and (211.4129,129.5159) .. (205.8964,208.2475) .. controls (250.9411,190.8018)
	and (304.2674,180.7163) .. (361.4277,180.7163) .. controls (417.0568,180.7163)
	and (469.0358,190.2795) .. (513.3027,206.8725) .. controls (507.1094,128.7941)
	and (440.6844,67.3100) .. (359.6464,67.3100) -- cycle;
	
	\fill[xshift=66,yshift=1,scale=0.0065,yscale=-1,xscale=-1,inner color=white,outer color=ce3dbdb!80!black,line width=0.2]
	(359.6346,72.6265) .. controls (315.4332,72.6265)
	and (275.6110,90.9215) .. (247.5096,120.2202) .. controls (283.3018,112.3912)
	and (322.3429,108.0640) .. (363.1971,108.0640) .. controls (400.7202,108.0640)
	and (436.6908,111.7186) .. (470.0408,118.3765) .. controls (442.0468,90.1452)
	and (402.9104,72.6265) .. (359.6346,72.6265) -- cycle;
	\end{tikzpicture}
}

\sectionfont{\LARGE}
\newcommand{\rpm}{\raisebox{.2ex}{$\scriptstyle\pm$}} % rpm is a smaller pm
\newcommand\tab[1][0.6cm]{\hspace*{#1}}
\newcommand\nl{\newline\tab}
\renewcommand\labelenumi{\textbf{\arabic{enumi}.}}
\renewcommand\labelenumii{\textbf{\theenumi.\arabic{enumii}.}}
\newcommand{\bigslant}[2]{{\raisebox{.2em}{$#1$}\left/\raisebox{-.2em}{$#2$}\right.}}
\newcommand{\suchthat}{\;\ifnum\currentgrouptype=16 \middle\fi|\;}

\renewcommand{\theenumii}{\textbf{.\arabic{enumii}.}}
%\renewcommand{\labelenumi}{\alph{enumi}}


\renewcommand\thesubsection{\arabic{subsection}.} 
\renewcommand\thesection{\arabic{section}} %Para que las secciones no se nombren como chapter.section (0.1, 0.2 etc)
\renewcommand\labelitemi{$\bullet$ }

\setcounter{secnumdepth}{3}

\linespread{1.2}

\title{Manifiesto de inutilidad}
\date{}
\begin{document}
	\maketitle
	\large{\textbf{\texttt{Holis.}}} \nl 
	Ayer me achicharré la espalda, y cada vez que salgo al sol mi piel me pide a gritos que me meta en casa. \nl Así que voy a hacerle caso y de paso intento poner mi cacao y mis bucles mentales en formato agradable, y así Yago tiene que joderse y leerlo.\nl \nl 
	Si $$\mathbb{S} = 
	\langle
	\begin{pmatrix}
	1 & 2 \\
	0 & 1 
	\end{pmatrix},
	\begin{pmatrix}
	1 & 0 \\
	2 & 1 
	\end{pmatrix}
	\rangle = \left\{ 
	\begin{pmatrix}
	a & b \\
	c & d 
	\end{pmatrix}
	\suchthat
	\begin{split}
	a \equiv d \equiv 1 &\mod 4 \\
	b \equiv c \equiv 0 &\mod 2 
	\end{split}
	\right\}, $$
	$\varphi$ es el homomorfismo sobreyectivo
	\begin{align*}	
	\varphi\colon SL_2(\mathbb{Z}) & \longrightarrow SL_2(\mathbb{Z}/ 4\mathbb{Z}) \\
			       \begin{pmatrix} 
					a & b\\
					c & d 
					\end{pmatrix} & \longmapsto
													\begin{pmatrix}
													a\bmod 4 & b \bmod 4 \\
													c\bmod 4 & d \bmod 4
													\end{pmatrix}
	\end{align*} y 
	$\phi$ es la función
	\begin{align*}	
	\phi\colon SL_2(\mathbb{Z}) & \longrightarrow SL_2(\mathbb{Z}) \\
	\begin{pmatrix} 
	a & b\\
	c & d 
	\end{pmatrix} & \longmapsto
	\begin{pmatrix}
	a\bmod 4 & b \bmod 2 \\
	c\bmod 2 & d \bmod 4
	\end{pmatrix}
	\end{align*}
	(Qué feo queda escrito lo de los módulos. D: ) \newline
	Por cierto, $\mathbb{S} = Ker(\phi)$
	\begin{itemize}
		\item \textbf{\texttt {Objetivo:}} Entender 
		$ \bigslant{SL_2(\mathbb{Z})}{\mathbb{S}}$ \nl
		\item \textbf{\texttt {Primer intento:}} Parecía que las clases de equivalencia de $ \bigslant{SL_2(\mathbb{Z})}{\mathbb{S}}$ tenían algo que ver, dada una matriz $\mathcal{M}$, con $\phi(\mathcal{M})$  (yo diría que dependen únicamente de $\phi(\mathcal{M})$, pero no con la boca muy llena... ) Por ejemplo, %$$
		
		\begin{align*}
		\phi(\mathcal{M}) = 
		\begin{pmatrix}
		m & 0 \\
		1 & n 
		\end{pmatrix} & \implies
		\phi(\overline{\mathcal{M}}) = 
		\phi(\mathcal{M}\mathbb{S}) = \left\{
		\begin{pmatrix}
		m & 0 \\
		1 & n 
		\end{pmatrix},
		\begin{pmatrix}
		m & 0 \\
		1 & n+2 
		\end{pmatrix},
		\right\}	\\			
		& \implies \overline{\mathcal{M}} \subseteq 
		 \phi^{-1}\left( 
		 \begin{pmatrix}
		 m & 0 \\
		 1 & n 
		 \end{pmatrix}
		  \right) 
		 \cup
		 \phi^{-1}\left( 
		 \begin{pmatrix}
		 m & 0 \\
		 1 & n+2 
		 \end{pmatrix}
		 \right)%$$
		 \end{align*}
		 
		 \tab Esta unión de conjuntos también ha quedado feísima.\nl 
		 ¿Es porque las $\phi$ son demasiado peques al lado de los paréntesis? Yo qué sé. Estética. \nl En fin.\nl 
		  El caso es que ese $\subseteq $ es un $=$, ya que 
		 $\text{para ninguna } \mathcal{N} \in \overline{\mathcal{N}} \neq \overline{\mathcal{M}} , \text{ se cumple que }\phi(\mathcal{N}) \in \phi(\overline{\mathcal{M}})$ \nl
		 E igual que hemos hecho con esta $\mathcal{M}$, para todos los posibles $\phi(\mathcal{N})$ podemos ver que $ \overline{\mathcal{N}}  = \phi^{-1}\left( \text{algún subconjunto de } SL_2(\mathbb{Z}/ 4\mathbb{Z}) \right)$ \nl\nl
		 \textbf{\texttt {Problema:}} No tiene por qué ser cierto que $
		 \begin{pmatrix}
		 a & b \\
		 c & d 
		 \end{pmatrix} \in SL_2(\mathbb{Z}/ 4\mathbb{Z}) \text{ para todos los } a,b,c,d \in \mathbb{Z}/ 4\mathbb{Z}
		 $ Abandonamos $\phi$ un ratito y usamos $\varphi$.
		 
		 \item \textbf{\texttt {Next step:}} Necesitamos saber quién es $SL_2(\mathbb{Z}/ 4\mathbb{Z})$, y una hermosa manera sería utilizando que $$SL_2(\mathbb{Z}/ 4\mathbb{Z}) = Im(\varphi) \cong \bigslant{SL_2(\mathbb{Z})}{Ker (\varphi)}$$ 
		 
		 Así que volvemos a tener el problema de estudiar un cociente. Pero, antes de eso, tenemos el problema de no tener ni puñetera idea de quién es $Ker(\varphi)$.\nl  Ojalá poder meter emojis en latex. A los emojis se les da bien describir sentimientos. Voy a ver si hay algún paquete para eso. \nl  
		 This is what I wanted. Mi cara al darme cuenta de que no tengo ni pajolera idea de quién es $Ker(\varphi)$ : \begin{center}
		 	\CreateStressedFace
		 \end{center}
	 	
	 	\item \textbf{\texttt {Cosas random antes de mirar el núcleo:}} $$ \varphi(\mathbb{S}) \cong \mathbb{Z}/ 2\mathbb{Z} \times \mathbb{Z}/ 2\mathbb{Z}, $$
	 	\tab Por otro lado, parece que $\frac{1}{4}$ de los elementos de $\mathbb{S}$ están en $Ker(\varphi)$. Que no sé si es $\frac{1}{4}$ de verdad, pero vaya. Aquellos donde los coeficientes $b$ y $c$ no son sólo pares, sino también múltiplos de 4. \nl Creo que esta observación es completamente inútil, pero me ha parecido curiosa.\nl  También me gustaría saber si sabiendo que $b$ y $c$ son pares, las matrices con $b$ y $c$ múltiplos de 4 son realmente  $\frac{1}{4}$ de $\mathbb{S}$. Estudiando índices hasta el $\infty$ y más allá.\nl
	 	Inciso: Lunes por la mañana, repasando esto antes de enviarlo. Acabo de decidir que no es que `` $\frac{1}{4}$ de los elementos de $\mathbb{S}$ estén en $Ker(\varphi)$ '', sino que $\frac{1}{4}$ de los elementos de $\mathbb{S}$ 	¡\textbf{son} $Ker(\varphi)$! Porque $$
	 	\varphi(\mathcal{M}) =  
	 	\begin{pmatrix}
	 	1 & 0 \\
	 	0 & 1 
	 	\end{pmatrix} \implies 
	 	\phi(\mathcal{M}) =  
	 	\begin{pmatrix}
	 	1 & 0 \\
	 	0 & 1 
	 	\end{pmatrix} \implies
	 	\mathcal{M} \in \mathbb{S}
	 	$$ 
	 	Por qué es como la quinta vez que me pongo a mirar $Ker(\varphi)$ y hasta ahora no veo algo tan tonto como eso, ¿eh? \begin{center} \CreateStressedFace \end{center} 
	 	\item \textbf{\texttt {Siguiente, un intento de ver el núcleo:}} 
	 	Imponiendo que $
	 	\mathcal{M} \in SL_2(\mathbb{Z}), \text{ y que } \varphi(\mathcal{M}) = 
	 	\begin{pmatrix}
	 	1 & 0 \\
	 	0 & 1 
	 	\end{pmatrix}
	 	$, lo más medianamente bonito que vi fue (después de hacer mal los numeritos quince mil veces, por lo que nada me asegura que esto esté bien) que $$
	 	Ker(\varphi) = \left\{ 
	 	\begin{pmatrix}
	 	4x+1 & 4y \\
	 	\frac{4xt + x + t}{4y} & 4t + 1 
	 	\end{pmatrix} 
	 	\suchthat 	 	
	 	\begin{split}
	 	x, t &\in \mathbb{Z}\\ y &\mid mcd(x, t)\\ 4 &\mid x + t
	 	\end{split}	 	
	 	\right\}
	 	$$
	 	\textbf{\texttt {Problema:}} Si $Ker(\varphi)$ ya es raro, ¡me puedo morir mirando el cociente! Lo intenté de manera fugaz, pero acababa encerrada en la misma línea de mirarlo todo en función de las clases de los coeficientes módulo 4... Y si no me había funcionado antes, no tenía mucha pinta de que fuera a funcionar ahora.
	 	
	 	
	 	\item \textbf{\texttt {Haciendo caso a Yago:}} En algún momento, Yago me dijo algo así como que podía conocer los elementos de $SL_2(\mathbb{Z}/4\mathbb{Z})$ usando el ordenador. \nl
	 	 A mi me parece que no, porque puedo sacar $\varphi$ de todas las matrices hasta cierta longitud de palabra, pero no tengo manera de saber si he terminado (si cojo matrices de mayor longitud, podría ser que $\varphi(\mathcal{M})$ fuera un nuevo elemento de $SL_2(\mathbb{Z}/4\mathbb{Z})$ que no me hubiera salido todavía). \nl 
	 	 Aun así, si $\mathcal{C}_n = \left\{ \text{matrices con longitud de palabra} \leq n\text{ según mis generadores }\right\}$	 	 
	 	 \begin{align*}
	 	  | \varphi(\mathcal{C}_1) | &= 4 \\	 	 
	 	  | \varphi(\mathcal{C}_2) | &= 15 \\	 	 
	 	  | \varphi(\mathcal{C}_3) | &= 31 \\	 	 
	 	  | \varphi(\mathcal{C}_4) | &= 43 \\	 	 
	 	  | \varphi(\mathcal{C}_5) | &= 47 \\	 	 
	 	  | \varphi(\mathcal{C}_6) | &= 48 \\	 	 
	 	  | \varphi(\mathcal{C}_7) | &= 48 \\	 	 
	 	  | \varphi(\mathcal{C}_8) | &= 48 \\	 	 
	 	  | \varphi(\mathcal{C}_9) | &= 48
	 	\end{align*}
	 	 \tab Por lo que tiene pinta de que $|SL_2(\mathbb{Z}/4\mathbb{Z})| = 48$, pero vaya. Que tampoco sé de qué me sirve eso si no puedo estar segura. 
	 	
	 	 
	 	 \item \textbf{\texttt {A otra cosa, mariposa:}} En una de mis desesperaciones de  ``Esto es terrible, no llego a ningún lado. Me aburro, así que voy a ver si aprendo algo guay leyendo cosas por internet'' encontré una dibujín muy bonito. Saltando de link en link, llegué \href{https://en.wikipedia.org/wiki/Modular_group}{aquí}. Evidentemente lo leo todo por encima porque si de paso me pongo a entenderlo, lo mismo me hago vieja. El caso es que daban una presentación de $PSL_2(\mathbb{Z})$... ¡¡con movimientos del señor Farey!! Y ver eso, con animación incluida, pues mola. Eso sí, ni puñetera idea de por qué esa era una presentación. Y justo después, hablan de trenzas. Y después, de tu plano. Y yo, leyendo medio dormida, me encuentro todas esas cosas juntas, y me emociono :)) \nl Por cierto, el tal $PSL_2(\mathbb{Z})$ está en tos laos. \nl
	\end{itemize}

	Espero haber transmitido adecuadamente mi inutilidad.\nl
	Help.\nl 
	Y así, queda consumido mi Domingo.\nl El día que escriba en \LaTeX,  la mitad de rápido que escribo en papel, me daré un premio.
	
\end{document}